% ===========================================================================
%  CAPITULO 1 — Como ler estas respostas
% ===========================================================================
\chapter{Como ler estas respostas}
\label{cap:como-ler}

Este documento responde, pergunta por pergunta, ao \emph{Questionário de
validação para os times} da Liga Digital (versão 1, 30/10/2026), escrito para
reunir as evidências que faltam ao Termo de Referência (TR) técnico e
pedagógico. As respostas foram levantadas no próprio repositório da
plataforma, e não em apresentações comerciais: o que está aqui como
implementado tem código, teste ou registro de decisão que o sustenta.

\section{Nomes do questionário e nomes da plataforma}

O questionário usa o vocabulário comercial. A plataforma usa outro, e as
evidências citadas estão escritas nele. A correspondência é esta:

\begin{center}
\footnotesize\renewcommand{\arraystretch}{1.3}
\begin{tabularx}{\linewidth}{@{}L{42mm}Y@{}}
\toprule
\textbf{No questionário} & \textbf{Na plataforma e nas evidências}\\
\midrule
Liga Digital & Plataforma Unooria: o hub web (professor, estudante, gestores
  e administração) e as experiências embutidas nele.\\
Missão & Experiência imersiva, identificada por um código como
  \cd{exp08ef1b} (8º ano do EF, 1º bimestre) ou \cd{exp02em1b} (2ª série do
  EM, 1º bimestre). Cada uma simula um aplicativo fictício (Unoovista,
  Unoochat, UnooriaAI, Instanooria\ldots).\\
Ficha da missão & A tela de detalhe da missão no painel do professor,
  alimentada pelo manifesto público da experiência
  (\arq{public/experiencia.json}).\\
Contextualização / personalização / setup & A criação da atividade pelo
  professor: tema e contexto escritos por ele, a partir dos quais a IA gera os
  conteúdos da simulação.\\
Aplicação para a turma & A \emph{liberação} da missão numa turma
  (\adrr{ADR-69}) somada à atividade que o professor põe no ar.\\
Tentativa do estudante & A \emph{rodada} ou sessão do estudante dentro de uma
  atividade.\\
Disponibilidade na rede & O interruptor do painel da Unooria
  (\adrr{ADR-87}): se o produto ou a missão existe naquele contrato, regional
  ou escola.\\
\bottomrule
\end{tabularx}
\end{center}

\section{Estrutura de cada resposta}

Cada pergunta ganha uma caixa com o código do questionário, um selo de status
e cinco campos, na ordem que o questionário pediu:

\begin{campos}
  \item[Hoje] o comportamento que existe, verificado no código.
  \item[Limitações] o que não funciona, só funciona em parte do catálogo, ou
    exige condição operacional.
  \item[Planejado] o que está em desenvolvimento ou decidido, com a fonte.
    Nada aqui é tratado como implementado.
  \item[Evidência] onde conferir: arquivo do repositório, ADR (Capítulo 18 da
    Documentação Técnica, \cd{PU-DOC-001}), história do Mapa de Produto
    (\cd{PU-PRD-001}) ou suíte e2e.
  \item[Lacuna] o que esta resposta não fecha e quem precisa fechar.
\end{campos}

Quando ajuda, a caixa é seguida de uma \textbf{proposta de redação para o TR}:
um requisito escrito como resultado observável com critério de aceite, sem
amarrar arquitetura, marca ou fornecedor, como pedem as observações finais do
questionário.

\section{Selos de status}

\begin{center}
\footnotesize\renewcommand{\arraystretch}{1.4}
\begin{tabularx}{\linewidth}{@{}L{48mm}Y@{}}
\toprule
\textbf{Selo} & \textbf{Significado}\\
\midrule
\implementado & Existe em todo o escopo perguntado, com evidência.\\
\parcial & Existe, mas só em parte do catálogo, sem algum dos controles
  pedidos ou com restrição operacional relevante.\\
\variapormissao & O comportamento difere entre missões; a resposta traz o
  recorte.\\
\naoimpl\newline\naoexiste & Não há código para isso hoje. Quando é
  compromisso da V1, a lacuna está descrita no Capítulo~\ref{cap:v1}.\\
\foradorepo & A resposta depende de material que não está no repositório
  (relatório de PoC, coleção didática, contrato de suporte) e precisa vir do
  time dono.\\
\bottomrule
\end{tabularx}
\end{center}

\begin{atencao}
\textbf{O repositório não responde tudo.} Perguntas sobre a prova de conceito
com o Instituto Paula Souza (P10), a coleção didática e o livro do professor
(E01, E02), a formação e o modelo de implantação (E03) e os compromissos
comerciais de suporte (parte de T09) dependem de documentos que não estão
versionados aqui. Para elas, o documento registra o que a plataforma oferece de
apoio e lista exatamente o que o time responsável precisa anexar.
\end{atencao}

\section{Data e base da verificação}

As respostas refletem a branch \cd{dev-local} em 30 de outubro de 2026. A
plataforma muda toda semana: antes de publicar o TR, confirme em ambiente de
homologação, com dados sintéticos, cada item marcado como implementado, como o
próprio questionário recomenda.
